% Test de charge de mod_ep : jeu de données d'une grande promotion, chronométrage des pages et
% des fonctions, corrections.
% Compilation : latexmk -pdf test-de-charge.tex
\documentclass[11pt,a4paper]{article}
\usepackage[french]{babel}
\usepackage[T1]{fontenc}
\usepackage[utf8]{inputenc}
\usepackage{lmodern,microtype,booktabs,tabularx,longtable,xcolor,hyperref,array}
\usepackage[margin=2.2cm]{geometry}
\hypersetup{colorlinks=true, linkcolor=blue!50!black, urlcolor=blue!50!black}
\newcommand{\code}[1]{\texttt{\detokenize{#1}}}
\newcolumntype{R}{>{\raggedleft\arraybackslash}p{1.75cm}}
\title{Test de charge de \code{mod_ep}}
\author{Sébastien Lefebvre}
\date{Novembre 2026 -- version 0.8.0 du module}

\begin{document}
\maketitle

\section{Méthode}

Le test porte sur un \textbf{vrai Moodle~4.5} (PostgreSQL~16, PHP~8.3), avec les quatre modules
\code{mod_ep}, \code{mod_epsynthesis}, \code{mod_stage} et \code{mod_stagesynthesis} installés,
servi par le serveur intégré de PHP à huit processus, sur une machine de quatre c\oe urs et
15~Go de mémoire. Le jeu de données (\code{docs/audit/banc/charge_setup.php}) représente une
promotion plusieurs fois plus grosse que celles d'une école vétérinaire :
\begin{itemize}
  \item un cours, une activité « Gestion des stages » (référents) et une activité EP ;
  \item 2\,000 étudiants répartis sur les six années, 150 enseignants (chacun référent d'environ
    treize étudiants et responsable d'un ou deux EP), 4 membres de la DEVE ;
  \item 200 EP au catalogue, chacun avec un syllabus publié de huit séances, un ou deux
    responsables, et pour un sur dix un prérequis ;
  \item 59\,550 crédits (30 par étudiant : inscriptions au catalogue, déclarations, ajouts de
    la DEVE) à tous les états -- 80~\% validés, 12~\% en attente à la première étape de leur
    circuit, 8~\% refusés -- avec l'historique de leurs décisions ;
  \item 6\,000 éléments du supplément au diplôme et 5\,000 décisions d'année.
\end{itemize}

Trois mesures :
\begin{enumerate}
  \item \textbf{Pages} (\code{pages.py}) : chaque page de l'activité, ouverte par un étudiant, un
    enseignant et un membre de la DEVE ; durée médiane de trois ouvertures, taille, nombre de
    lectures et d'écritures en base (pied de page de Moodle avec \code{perfdebug}), erreurs.
  \item \textbf{Fonctions} (\code{fonctions.php}) : fonctions de lecture de \code{locallib.php},
    médiane de cinq appels.
  \item \textbf{Traitements lourds} : export complet, restauration de cet export dans une autre
    activité, durée et mémoire.
\end{enumerate}
Les scripts sont dans \code{docs/audit/banc/} et rejouables (\code{README.md}).

\section{Problèmes trouvés et corrections}

La première mesure a relevé six points, tous corrigés dans la version~0.8.0 :

\begin{longtable}{@{}p{4.2cm}p{5.3cm}p{6.2cm}@{}}
\toprule
\textbf{Page ou traitement} & \textbf{Constat} & \textbf{Correction} \\
\midrule\endhead
Validation (DEVE) & 3,8~Mo et 0,8~s : la file « En attente de votre décision » affichait les
  7\,000 demandes en attente, sans pagination. & File paginée en base (40 par page), avec son
  total. \\
Éléments du supplément (DEVE) & 2\,083 requêtes, 1,6~s, 2,4~Mo : 2\,000 éléments en attente
  affichés d'un coup, preuves lues élément par élément, liste complète paginée en PHP après
  lecture de tout. & File et liste paginées en base ; preuves de la page lues en une requête. \\
Catalogue (étudiant) & 315 requêtes : prérequis manquants calculés EP par EP, deux fois par
  ligne. & Calcul groupé pour tout le catalogue (4 requêtes). \\
Participants d'un EP & Prérequis recalculés étudiant par étudiant. & Calcul groupé pour tous les
  inscrits en attente. \\
Export complet & 34~s, plus de 1~Go de mémoire (épuisement en ligne de commande), 66\,500
  requêtes : classeur PhpSpreadsheet tenu en mémoire (800\,000 cellules), fichiers lus crédit
  par crédit. & Classeur écrit en flux (OpenSpout, fourni avec Moodle), données JSON écrites en
  un passage sur les crédits lus par paquets, fichiers lus par zone, bilans par paquets de 500
  étudiants : \textbf{10~s, 102~Mo, 970 requêtes}. \\
Restauration de l'export & 7~min sous verrou, 198\,000 requêtes : un contrôle de doublon par
  ligne, dans une table qui grossit à chaque insertion, devenait de plus en plus lent. & Index
  des enregistrements présents lu une fois, décisions insérées par lots : \textbf{15~s,
  70\,800 requêtes} pour 59\,550 crédits. \\
\bottomrule
\end{longtable}

\section{Pages, après correction}

Durées médianes côté serveur, une seule personne connectée. Les pages HTML ne dépassent pas
0,2~s ; le nombre de requêtes ne dépend plus de la taille de la promotion, sauf pour l'émission
des suppléments, dont chaque ligne rassemble le supplément de l'étudiant (20 étudiants par page,
voir plus bas).

\begin{longtable}{@{}p{1.9cm}p{7.6cm}RRR@{}}
\toprule
\textbf{Rôle} & \textbf{Page} & \textbf{Durée (s)} & \textbf{Requêtes} & \textbf{Taille (Ko)} \\
\midrule\endhead
Étudiant & Tableau de bord (\code{view.php}) & 0,11 & 41 & 175 \\
 & Catalogue (200 EP) & 0,11 & 44 & 330 \\
 & Déclarer un EP & 0,11 & 64 & 216 \\
 & Mon supplément au diplôme & 0,08 & 55 & 160 \\
 & Fiche du syllabus & 0,06 & 40 & 159 \\
 & Fiche du syllabus en PDF & 0,10 & -- & 146 \\
Enseignant & Tableau de bord & 0,12 & 45 & 163 \\
 & Validation & 0,09 & 47 & 210 \\
 & Pilotage & 0,08 & 45 & 163 \\
 & Propositions & 0,06 & 39 & 157 \\
 & Éléments du supplément & 0,07 & 43 & 170 \\
 & Participants d'un EP & 0,12 & 51 & 338 \\
DEVE & Tableau de bord & 0,12 & 46 & 184 \\
 & Pilotage (2\,000 étudiants, paginé) & 0,10 & 46 & 184 \\
 & Pilotage, recherche, page 5 & 0,09 & 46 & 184 \\
 & Situation d'un étudiant & 0,07 & 48 & 184 \\
 & Validation & 0,12 & 46 & 217 \\
 & Validation, filtre « en attente », page 3 & 0,10 & 46 & 218 \\
 & Catalogue & 0,10 & 44 & 382 \\
 & Gestion du catalogue & 0,10 & 44 & 426 \\
 & Participants d'un EP & 0,13 & 48 & 321 \\
 & Types, minimums, administration & 0,06--0,07 & 40--42 & 163--169 \\
 & Validation des années & 0,09 & 48 & 216 \\
 & Éléments du supplément & 0,10 & 48 & 235 \\
 & Émission des suppléments & 0,18 & 363 & 218 \\
 & Supplément d'un étudiant, aperçu PDF & 0,09 / 0,13 & 59 / -- & 168 / 152 \\
 & Réglages du supplément & 0,10 & 71 & 267 \\
 & Rentrée, transfert du paramétrage, export complet, import du catalogue & 0,08--0,09 & 44--64 &
   164--204 \\
 & Export CSV par étudiant & 0,70 & -- & 186 \\
 & Export CSV détaillé (59\,550 lignes) & 3,58 & -- & 9\,738 \\
 & Livret PDF du catalogue (200 EP) : production, puis livret conservé & 4,86 / 0,03 & -- & 1\,320 \\
 & Export complet (ZIP, 60\,000 crédits) & 9,49 & -- & 4\,800 \\
 & Paramétrage JSON avec catalogue & 0,15 & -- & 1\,425 \\
\bottomrule
\end{longtable}

\section{Fonctions de lecture}

\begin{longtable}{@{}p{10cm} R R@{}}
\toprule
\textbf{Fonction} & \textbf{ms} & \textbf{Requêtes} \\
\midrule\endhead
Demandes en attente de la DEVE : total & 7,4 & 1 \\
Demandes en attente de la DEVE : une page de 40 & 8,7 & 1 \\
Demandes en attente d'un enseignant (toutes) & 1,0 & 1 \\
Liste filtrée avec recherche sur le nom, page 3 & 24,2 & 1 \\
Bilan d'une page de pilotage (40 étudiants) & 11,4 & 5 \\
Bilan de 500 étudiants & 114,3 & 5 \\
Bilan d'un étudiant & 0,4 & 3 \\
Places prises et statuts de tous les EP & 8--9 & 1 \\
Prérequis manquants d'un étudiant, catalogue entier & 1,0 & 4 \\
Contenu du supplément d'un étudiant & 2,3 & 13 \\
Éléments du supplément en attente, page de 40 & 2,6 & 1 \\
Décisions d'année de 500 étudiants & 3,1 & 1 \\
Export du paramétrage avec catalogue & 86,7 & 812 \\
\bottomrule
\end{longtable}

\section{Points restants}

\begin{itemize}
  \item \textbf{Émission des suppléments} : environ 18 requêtes par étudiant de la page (le
    contenu complet du supplément est rassemblé pour signaler ce qui manque), soit 360 pour une
    page de 20 et 0,18~s : borné par la pagination, sans effet perceptible.
  \item \textbf{Export du paramétrage} : quatre requêtes par EP du catalogue (prérequis, syllabus,
    séances, intervenants), soit 0,09~s pour 200 EP. Opération rare.
  \item \textbf{Restauration} : le fichier \code{donnees.json} est lu d'un bloc (300~Mo de
    mémoire pour 60\,000 crédits). La page relève la limite de mémoire de PHP pour cette
    opération ; au-delà de 150\,000 crédits dans une même activité, mieux vaut restaurer en
    ligne de commande (\code{docs/audit/concurrence/worker.php}).
  \item \textbf{Livret PDF} : 4,9~s pour produire le livret de 200 EP de huit pages de syllabus
    chacun ; il est ensuite conservé et resservi en 0,03~s tant que le catalogue, les syllabus, les
    responsables et les réglages des PDF ne changent pas.
\end{itemize}

\section{Conclusion}

Avec 2\,000 étudiants, 200 EP et 60\,000 crédits dans une seule activité, toutes les pages
s'ouvrent en moins de 0,2~s côté serveur, avec un nombre de requêtes qui ne dépend pas de la
taille de la promotion. Les traitements de masse restent dans des durées acceptables pour une
requête web (export complet 10~s, restauration 15~s, livret 5~s la première fois) et une mémoire bornée. La
tenue sous accès simultanés fait l'objet du rapport \emph{Test de concurrence}.
\end{document}
